%\documentclass[12pt]{scrreprt}
\documentclass[12pt]{report} 

% language may be romanian or english (default is english)
% type may be bachelor or master (default is bachelor)
\usepackage[language=english, type=bachelor]{style}
\usepackage{amsmath}
\usepackage{amssymb}

% Optional figures: place PNG/PDF under ./figures/ (see THESIS_FIGURES.md).
% Compile from the directory that contains the figures/ folder (project root).
% If the file is missing, no \includegraphics runs (avoids pdftex.def ``not found'' errors).
\newcommand{\thesisfig}[2]{%
  \IfFileExists{#2}{%
    \includegraphics[width=#1]{#2}%
  }{%
    \fbox{\parbox{#1}{\centering\vspace{1.1cm}\footnotesize\textit{Add file:} \texttt{\detokenize{#2}}\vspace{1.1cm}}}%
  }%
}

%\geometry{a4paper,top=2.5cm,left=3cm,right=2.5cm,bottom=2.5cm}
%in style
%controlling the appearance of your headers and footers
\usepackage{fancyhdr}
\setlength{\headheight}{27.35799pt}
\pagestyle{fancy}
\lhead{}
\chead{}
\renewcommand{\headrulewidth}{0.2pt}
\renewcommand{\footrulewidth}{0.2pt}

\begin{document}

\specialization{Computer Science}	
\title{Augmented Reality-Based Interior Design: A Mobile Application for Parametric Room Scanning and Interactive 3D Scene Reconstruction}					   
\author{Mara Ielciu}											
\supervisor{Professor Dr. Darabant Sergiu Adrian}				
				
\maketitle


\newpage
\thispagestyle{empty}
\mbox{}
\newpage
\pagenumbering{roman} 

\cleardoublepage
ABSTRACT
\vspace{0.5cm}	
\hrule
\vspace{0.5cm}	
%\cleardoublepage

This bachelor thesis presents an iOS application for augmented reality--based interior design. Users scan a real room with Apple's RoomPlan API to obtain a parametric, semantically labelled model (\texttt{CapturedRoom}: walls, doors, windows, and detected furniture with metric dimensions). The scan is loaded into a non-AR three-dimensional editor built with RealityKit, where detected furniture can be replaced with USDZ models from a library or produced by Object Capture's on-device photogrammetry pipeline. Edits persist locally (JSON scan data under the app Documents directory, Codable room models). The project also supports importing pretrained 3D Gaussian Splatting assets (\texttt{.ply}) and viewing them with a custom Metal renderer, enabling visual comparison with mesh-based reconstruction.

Chapter~2 reviews photogrammetry (structure from motion, multi-view stereo, Poisson surface reconstruction), depth fusion via truncated signed distance functions, and how Object Capture versus RoomPlan map onto image-based versus depth-assisted reconstruction. Chapter~3 summarises 3D Gaussian Splatting: Gaussian primitives, projection and splatting, differentiable optimisation, and why meshes remain preferable for interactive furniture replacement. Chapters~4 and~5 state requirements and describe the software architecture (SwiftUI, MVVM-oriented structure, major views and services), data flow, and testing. The main contribution is an integrated, on-device workflow from room scan to persistent, editable virtual staging, with an optional path to plenoptic visualisation via 3DGS imports.



\tableofcontents


\newpage
\pagenumbering{arabic}

% ---------------------------------------------------------------------------
% Context and related work (3 paragraphs + 3 bibliographic resources)
% ---------------------------------------------------------------------------
\section*{Context and Related Work}
\addcontentsline{toc}{section}{Context and Related Work}

Indoor scene understanding and 3D reconstruction have been central to computer vision research, with datasets and benchmarks driving progress in semantic segmentation and object detection. ScanNet \cite{dai2017scannet} is a large-scale dataset of richly annotated 3D reconstructions of indoor scenes, with 2.5 million views across more than 1,500 scans and instance-level semantic segmentations. That work established the importance of combining geometric reconstruction with semantic labelling (walls, furniture, openings) for applications such as robotics and spatial reasoning. The present application aligns with this direction by leveraging Apple's RoomPlan to produce a parametric, semantically annotated room model from a consumer device, rather than from offline processing of RGB-D videos.

The use of augmented reality for interior design and furniture visualisation has been explored in both industry and academia. K\'an et al.\ \cite{kan2021ar} present an AR system for automatic interior design based on a hierarchical tree of procedural rules, enabling interactive and personalised furniture configuration with a recommender component; they compare AR-based visualisation to traditional methods and report positive user feedback. Such studies confirm that AR can support decision-making and spatial reasoning when users need to envision furniture and finishes in their own spaces. This thesis extends that idea by allowing users not only to place virtual furniture, but to \emph{replace} detected real furniture in a reconstructed 3D scene, thus supporting ``what-if'' redesign directly tied to the current room layout.

User experience in mobile AR interior design applications has been studied to improve usability and adoption. Kandil et al.\ \cite{kandil2021ux} evaluate interior design mobile AR applications using established scales such as the System Usability Scale (SUS) and the Handheld Augmented Reality Usability Scale (HARUS), with a focus on gesture-based interaction and one-handed use. Their findings underline the need for clear onboarding, stable tracking, and intuitive manipulation of virtual objects. The application developed in this thesis addresses these concerns by using RoomPlan's guided scanning flow, RealityKit's responsive rendering, and SwiftUI's declarative interface to keep interactions consistent and learnable.


% ═══════════════════════════════════════════════════════════════════════════
%  CHAPTER 1 — INTRODUCTION
% ═══════════════════════════════════════════════════════════════════════════

\chapter{Introduction}
\label{ch:introduction}

This thesis presents the design and development of a mobile application that combines \textbf{parametric room scanning} and \textbf{interactive 3D scene reconstruction} for augmented reality (AR)--based interior design. The goal is to allow users to scan a room with a consumer device, obtain a structured representation of the space---including walls, openings, and furniture---and then visualise and edit the scene by replacing detected real furniture with virtual alternatives. The application thus supports ``what-if'' redesign directly tied to the user's own environment.

\section{Problem Statement and Goals}

Interior design decisions are often difficult to visualise without seeing furniture and finishes in place. Traditional tools range from 2D floor plans and mood boards to physical samples; these do not provide an immersive, spatially accurate preview of the final result. AR has emerged as a way to overlay virtual content onto the real world, but many solutions require the user to remain in the physical space with the device pointed at the room. A persistent, editable 3D model of the room---obtained once by scanning---enables users to explore design options at any time, without being on-site.

The main goals of this work are:

\begin{enumerate}
    \item To integrate \textbf{room-scale} capture (Apple's RoomPlan) to obtain a parametric, semantically annotated model of the room, including detected furniture with poses and dimensions.
    \item To provide an \textbf{interactive 3D room editor} where users can replace detected furniture with virtual 3D models (USDZ), with gesture-based manipulation and persistence.
    \item To support \textbf{object-scale} capture (Apple's Object Capture) for scanning individual objects into 3D models, and to allow \textbf{comparison with 3D Gaussian Splatting} reconstructions as an alternative representation.
\end{enumerate}

\section{Scope of the Thesis}

The thesis focuses on the theoretical foundations and the design of the application. \textbf{Chapter~\ref{ch:theory1}} covers the motivation for the work, the principles of photogrammetry, and the role of Object Capture and Room Plan in the pipeline. \textbf{Chapter~\ref{ch:theory2}} introduces 3D Gaussian Splatting as an alternative to mesh-based reconstruction and explains how it is used in the application for viewing and comparison. Later chapters address the implementation, architecture, and conclusions.

\section{Document Structure}

\begin{itemize}
    \item \textbf{Chapter~\ref{ch:theory1}} --- Theoretical Foundations I: motivation, photogrammetry, Object Capture, and Room Plan.
    \item \textbf{Chapter~\ref{ch:theory2}} --- Theoretical Foundations II: 3D Gaussian Splatting.
    \item \textbf{Chapter~4} --- Implementation and application architecture.
    \item \textbf{Chapter~6} --- Conclusions and future work.
\end{itemize}

%\addcontentsline{toc}{chapter}{Introducere}
%\addcontentsline{toc}{chapter}{Introduction}

\chapter{Photogrammetry, Depth Fusion, and Apple Capture APIs}
\label{ch:theory1}

This chapter explains the geometric and algorithmic ideas behind the capture pipelines used in the thesis application. It first reviews photogrammetry and the classical chain structure from motion (SfM) $\rightarrow$ multi-view stereo (MVS) $\rightarrow$ Poisson surface reconstruction \cite{luhmann2014photogrammetry,lowe2004sift,furukawa2010pmvs,kazhdan2006poisson}. It then describes depth fusion with a truncated signed distance function (TSDF) as used in real-time RGB-D and LiDAR systems \cite{newcombe2011kinectfusion}, which underpins ARKit-style scene meshes that RoomPlan builds upon. A comparison between image-only object reconstruction and depth-assisted room understanding clarifies why Object Capture and RoomPlan behave differently. Finally, the chapter maps these ideas to Apple's Object Capture and RoomPlan APIs and documents the 'CapturedRoom' data model.


\section{Motivation}
\label{sec:motivation}

Users redesign interiors more confidently when they can see alternatives in context. Traditional representations (2D plans, catalog photos, or material swatches) omit depth, occlusion, and the relation between existing objects and new ones. Augmented reality can composite virtual furniture into a live video stream, but that mode ties the experience to being physically present and holding the device. If the room is first captured as a structured 3D model, the user can reopen and edit that model later: the scan becomes a persistent spatial substrate for experimentation.

Research datasets such as ScanNet \cite{dai2017scannet} demonstrated that indoor applications benefit from joint geometry and semantics (which surfaces are walls, which objects are chairs or tables). RoomPlan brings a consumer subset of that idea to supported iPhones and iPads: it outputs not only geometry but labelled entities with rough dimensions. This thesis uses that output as the basis for furniture replacement, substituting USDZ models for detected instances, which goes beyond merely placing a free-floating object on a manually identified floor plane.

\section{Photogrammetry and the Classical Reconstruction Pipeline}
\label{sec:photogrammetry}

\subsection{Definition and triangulation}

\textbf{Photogrammetry} is the discipline of extracting metric information from photographs \cite{luhmann2014photogrammetry}. In 3D reconstruction, the elementary geometric operation is 
triangulation: corresponding points in two or more calibrated images determine a 3D point by intersecting rays. Stable triangulation requires baseline between cameras and sufficient overlap so that the same surface points appear in multiple views.
\begin{figure}[htbp]
\centering
\input{tikz_diagrams/ch2_photogrammetry_concept.tex}
\caption{Triangulation and Bundle Adjustment: corresponding image points from calibrated cameras determine a 3D point~$\mathbf{X}$; Bundle Adjustment jointly minimises reprojection error over all camera poses~$\mathbf{P}_j$ and scene points~$\mathbf{X}_i$.}
\label{fig:ch2_photo_concept}
\end{figure}

\subsection{From features to textured mesh}

A widely used multi-view pipeline proceeds in stages \cite{luhmann2014photogrammetry,lowe2004sift,furukawa2010pmvs,kazhdan2006poisson}:

\begin{enumerate}
    \item \textbf{Feature detection and matching}: Distinctive keypoints (e.g.\ SIFT \cite{lowe2004sift}) are extracted and matched across images.

    \item \textbf{Structure from Motion}: Camera intrinsics and extrinsics are estimated together with a \emph{sparse} 3D point cloud. \textbf{Bundle adjustment} refines parameters by minimising reprojection error across all observations. Symbolically, one minimises $\sum_{i,j} \|\mathbf{x}_{ij} - \pi(\mathbf{P}_j, \mathbf{X}_i)\|^2$ over projection matrices $\mathbf{P}_j$ and points $\mathbf{X}_i$, where $\mathbf{x}_{ij}$ are measured image points and $\pi$ is the projection function.

    \item \textbf{Multi-view stereo}: Dense correspondence and depth estimation across views yield a dense point cloud or per-view depth maps \cite{furukawa2010pmvs}.

    \item \textbf{Surface reconstruction}: The dense samples are converted to a watertight or oriented surface, commonly using \textbf{Poisson surface reconstruction} \cite{kazhdan2006poisson}, which solves for an indicator function whose gradient aligns with surface normals, then extracts a mesh (e.g.\ Marching Cubes).

    \item \textbf{Texture mapping}: The mesh is parameterised; colours are blended from source images onto the surface.
\end{enumerate}
\begin{figure}[htbp]
\centering
\input{tikz_diagrams/fig2_2_sfm_mvs}
\caption{SfM--MVS reconstruction pipeline: from multi-view image input
through feature detection, Structure from Motion (SfM), dense Multi-View
Stereo (MVS), to surface reconstruction and textured mesh output. Pipeline
stages follow the standard formulation in~\cite{luhmann2014photogrammetry,lowe2004sift,furukawa2010pmvs,kazhdan2006poisson}.}
\label{fig:ch2_photo_pipeline}
\end{figure}

This end-to-end path is what Section~\ref{sec:on_device_pipeline} relates to \textbf{Object Capture}.

\section{Depth Fusion and TSDF-Based Scene Meshing}
\label{sec:depth_fusion}

Room-scale tracking on mobile devices fuses depth (LiDAR or stereo) with visual-inertial odometry. A standard volumetric representation is the TSDF: for each voxel in a bounded volume, one stores a (truncated) signed distance to the nearest surface along with a weight; each new depth frame updates voxels along viewing rays \cite{newcombe2011kinectfusion}. Over time, noise averages down and a coherent surface emerges. The mesh shown in AR debugging views is typically an isosurface of that field.

This reconstruction is driven by measured depth: textureless walls still contribute if the depth sensor returns reliable values. By contrast, MVS in photogrammetry relies on photometric consistency and can fail on untextured or specular surfaces. Semantic interpretation (which planar patch is a ``wall'' and which blob is a ``sofa'') is not solved by the TSDF alone; RoomPlan adds learning-based detection and classification on top of ARKit scene understanding.

\section{Image-Only versus Depth-Assisted Reconstruction}
\label{sec:depth_vs_photo}

\textbf{Object Capture} is primarily an image-driven pipeline (with poses from ARKit): geometry is inferred from many RGB views through stages analogous to SfM, MVS, and Poisson (or similar) meshing, followed by texturing. \textbf{Room Plan} is depth- and semantics-first: it targets a structured, parametric room description ('CapturedRoom') rather than a single closed mesh of the whole interior. Table~\ref{tab:depth_vs_photo} contrasts the two paradigms at the level relevant to this thesis.

\begin{table}[htbp]
\centering
\caption{Comparison of image-centric object reconstruction (Object Capture) and depth-assisted parametric room capture (Room Plan). The Room Plan branch builds on real-time depth integration in the TSDF family \cite{newcombe2011kinectfusion}, augmented by learned structure and semantics.}
\label{tab:depth_vs_photo}
\begin{tabular}{|p{0.28\textwidth}|p{0.30\textwidth}|p{0.30\textwidth}|}
\hline
\textbf{Aspect} & \textbf{Object Capture} & \textbf{Room Plan} \\
\hline
Primary sensing & Dense RGB coverage & LiDAR depth + RGB + IMU \\
\hline
Dominant geometry cue & Photo-consistency (MVS) & TSDF fusion + plane fitting \\
\hline
Typical output & Textured triangle mesh (USDZ) & Parametric walls, openings, objects \\
\hline
Semantics & Optional masking & Built-in categories (e.g.\ chair, table) \\
\hline
\end{tabular}
\end{table}

\section{On-Device Capture Pipeline}
\label{sec:on_device_pipeline}

Figure~\ref{fig:capture_branches} compares two parallel pipelines.  
The left column is the photogrammetric/Object Capture path:
multi-view RGB images and poses $\rightarrow$ feature matching $\rightarrow$ SfM
(sparse structure) $\rightarrow$ MVS (dense geometry) $\rightarrow$ surface reconstruction
and texturing $\rightarrow$ USDZ output.
The right column is the RoomPlan path:
LiDAR/RGB/IMU input $\rightarrow$ VIO tracking $\rightarrow$ TSDF/depth fusion and mesh anchors
$\rightarrow$ plane/structure inference $\rightarrow$ semantic object/opening detection
$\rightarrow$ parametric fitting $\rightarrow$ 'CapturedRoom' output.
Apple’s internal implementations are proprietary; therefore, the figure
is a conceptual mapping of stages to public APIs.

% \begin{table}[htbp]
% \centering
% \caption{Conceptual mapping from classical stages to Apple's on-device APIs.}
% \label{tab:pipeline_to_api}
% \begin{tabular}{|l|p{0.35\textwidth}|p{0.35\textwidth}|}
% \hline
% \textbf{Stage} & \textbf{Object Capture} & \textbf{Room Plan} \\
% \hline
% Camera trajectory & ARKit VIO; session stores images & ARKit world tracking during scan \\
% \hline
% Sparse / dense geometry & Photogrammetry engine (SfM/MVS class) & TSDF-class fusion + parametric fitting \\
% \hline
% Surface model & Textured mesh & Structural elements + object boxes \\
% \hline
% Semantics & Background masking (when enabled) & Wall / door / window / object classes \\
% \hline
% Serialised output & USDZ file & 'CapturedRoom' (Codable) \\
% \hline
% \end{tabular}
% \end{table}

\begin{figure}[H]
\centering
\input{tikz_diagrams/fig2_3_dual_pipeline}
% \thesisfig{1\textwidth}{figures/fig2_3_dual_pipeline.pdf}
\caption{Parallel capture pipelines: Object Capture follows an 
image-based reconstruction chain (left column); RoomPlan combines 
LiDAR depth fusion with parametric and semantic inference to produce 
a 'CapturedRoom' (right column).}
\label{fig:capture_branches}
\end{figure}

\subsection{Object Capture}

ObjectCaptureSession guides the user around an object so that views overlap; PhotogrammetrySession reconstructs a mesh and materials, outputting USDZ. That format is well supported in RealityKit and matches the furniture assets used in the editor.

\subsection{Room Plan and CapturedRoom}

Room Plan exposes types such as walls, doors, windows, and CapturedRoom.Object instances with 'category', 'dimensions', and 'transform'. Apple continues to extend taxonomy and behaviour \cite{apple2024roomplan2}. Because 'CapturedRoom' conforms to 'Codable', applications can persist scans as JSON or similar, as done in the implementation (Chapter~\ref{ch:implementation}) under a per-room file in the Documents hierarchy.

\subsection{Feature-to-API mapping}

Table~\ref{tab:feature_api} ties user-visible capabilities to frameworks.

\begin{table}[h]
\centering
\caption{Major application features and the APIs or artefacts that implement them.}
\label{tab:feature_api}
\begin{tabular}{|l|l|p{0.42\textwidth}|}
\hline
\textbf{Feature} & \textbf{API / artefact} & \textbf{Role} \\
\hline
Room scanning & RoomPlan & Produce 'CapturedRoom' \\
\hline
Room persistence & 'Codable' + file I/O & Reload scans after app restart \\
\hline
3D editor & RealityKit (non-AR) & Visualise and edit room \\
\hline
Object scanning & Object Capture & Author USDZ furniture \\
\hline
Splat viewing & Metal + '.ply' & Compare with meshes \\
\hline
\end{tabular}
\end{table}

\section{Limitations}
\label{sec:limitations_capture}

Photogrammetry remains sensitive to glass, mirrors, and uniform surfaces. Room Plan can miss small objects, merge adjacent pieces, or mislabel categories; coarse bounding boxes are not CAD precision. These issues motivate manual touch-up in the editor and, for research-oriented quality, optional export to Blender for refinement; parametric RoomPlan geometry is exported as glTF boxes and markers, with USDZ furniture applied via post-import script, including 3DGS training pipelines discussed in Chapter~\ref{ch:theory2}.  


% ═══════════════════════════════════════════════════════════════════════════
%  CHAPTER 3 — THEORETICAL FOUNDATIONS II: 3D GAUSSIAN SPLATTING
% ═══════════════════════════════════════════════════════════════════════════

\chapter{Theoretical Foundations II: 3D Gaussian Splatting}
\label{ch:theory2}

Object Capture and the photogrammetry pipeline of Chapter~\ref{ch:theory1} produce a \textbf{mesh-based} representation: a triangle mesh with textures. This chapter introduces \textbf{3D Gaussian Splatting} (3DGS) \cite{kerbl20233dgs}, an alternative scene representation that uses a set of 3D Gaussian primitives splatted onto the image plane and alpha-blended for rendering. 3DGS yields high visual quality and real-time rendering, and is differentiable, so it can be trained from multi-view images without explicit mesh or surface extraction. The application supports both mesh-based (Object Capture) and 3DGS reconstructions, allowing users to compare the two representations.

% ---------------------------------------------------------------------------
\section{What is 3D Gaussian Splatting}
\label{sec:3dgs_intro}

3D Gaussian Splatting represents a scene as a large set of \textbf{3D Gaussian primitives}. Each primitive has a position in 3D space, an anisotropic shape (scale and orientation), an opacity, and a colour (or view-dependent colour). To render an image, each Gaussian is \textbf{projected} onto the 2D image plane, where it becomes a 2D Gaussian (or an elliptical blob). The primitives are then \textbf{sorted by depth} and \textbf{alpha-blended} back-to-front to produce the final pixel colours. The entire process is differentiable with respect to the Gaussian parameters, so that gradients from an image loss can be backpropagated during training to optimise positions, shapes, opacity, and colour from multi-view photographs.

Unlike mesh-based reconstruction, 3DGS does not require an explicit surface (e.g.\ Poisson reconstruction or mesh extraction). The scene is represented implicitly by the set of Gaussians, which can model complex geometry and view-dependent effects. In the application developed in this thesis, 3DGS models are typically trained off-device (e.g.\ from images exported by Object Capture or from a separate capture); the app then \emph{imports} the resulting \texttt{.ply} file and renders it with a custom Metal-based splatting pipeline for viewing and comparison with Object Capture meshes.

% ---------------------------------------------------------------------------
\section{Representation: 3D Gaussians}
\label{sec:3dgs_representation}

Each primitive is an \textbf{anisotropic 3D Gaussian} defined by:

\begin{itemize}
    \item \textbf{Position} $\boldsymbol{\mu} \in \mathbb{R}^3$: the centre of the Gaussian in world space.

    \item \textbf{Covariance} $\boldsymbol{\Sigma} \in \mathbb{R}^{3 \times 3}$: a symmetric positive definite matrix that defines the shape and orientation of the Gaussian. For stability and differentiability, 3DGS parameterises it using \textbf{scale} $\mathbf{s} \in \mathbb{R}^3$ (positive per-axis scales) and \textbf{rotation} as a unit quaternion $\mathbf{q}$. Then $\boldsymbol{\Sigma} = \mathbf{R} \mathbf{S} \mathbf{S}^\top \mathbf{R}^\top$, where $\mathbf{R}$ is the rotation matrix of $\mathbf{q}$ and $\mathbf{S} = \operatorname{diag}(\mathbf{s})$.

    \item \textbf{Opacity} $\alpha \in [0,1]$: used in alpha-blending during rasterisation.

    \item \textbf{Colour}: typically stored as spherical harmonics (SH) coefficients for view-dependent effects, or as a single RGB value per primitive. The application's \texttt{.ply} parser and Metal renderer use per-Gaussian RGB (linear space).
\end{itemize}

The 3D Gaussian function is
\[
G(\mathbf{x}) = \exp\bigl(-\tfrac{1}{2}(\mathbf{x}-\boldsymbol{\mu})^\top \boldsymbol{\Sigma}^{-1} (\mathbf{x}-\boldsymbol{\mu})\bigr).
\]
For rendering, this 3D function is projected to 2D in screen space, as described in the next section.

% ---------------------------------------------------------------------------
\section{Splatting: Projection and Rasterisation}
\label{sec:3dgs_splatting}

Rendering proceeds by \textbf{splatting} each Gaussian onto the image plane:

\begin{enumerate}
    \item \textbf{View projection}: The 3D mean $\boldsymbol{\mu}$ is projected to 2D screen coordinates using the camera's view and projection matrices. The 3D covariance $\boldsymbol{\Sigma}$ is projected to a 2D covariance $\boldsymbol{\Sigma}'$ in screen space using the Jacobian of the projection at $\boldsymbol{\mu}$ \cite{kerbl20233dgs}. This 2D covariance defines an ellipse in the image plane.

    \item \textbf{Sorting}: Primitives must be composited in a consistent order for alpha-blending. Typically they are sorted by depth (e.g.\ by the $z$-coordinate of $\boldsymbol{\mu}$ in camera space) so that back-to-front blending is well-defined. Tile-based sorting can also be used for efficiency.

    \item \textbf{Rasterisation}: Each Gaussian is drawn as a quad or ellipse. The fragment shader evaluates the 2D Gaussian weight (or an approximation thereof), multiplies by opacity and colour, and composites with the frame buffer using alpha blending: $\mathbf{c}_{\mathrm{out}} = \alpha \cdot \mathbf{c}_{\mathrm{G}} + (1-\alpha) \mathbf{c}_{\mathrm{in}}$.
\end{enumerate}

Because each step is differentiable, the full pipeline allows gradients from a photometric loss (e.g.\ L1 or L2 between rendered and ground-truth pixels) to be backpropagated to update the Gaussian parameters during training. Training is usually performed off-device: given multi-view RGB images and camera poses (e.g.\ from COLMAP or from an AR/Object Capture session), the Gaussians are initialised from a sparse point cloud, then optimised by minimising the rendering loss, with optional pruning and densification to control the number of primitives. The output is a \texttt{.ply} file containing per-vertex (per-Gaussian) attributes: position, colour, opacity, scale, and rotation.

% ---------------------------------------------------------------------------
\section{Relation to Mesh-Based Reconstruction}
\label{sec:3dgs_vs_mesh}

In the context of this thesis, the two main scene representations are:

\begin{itemize}
    \item \textbf{Object Capture (mesh)}: Explicit geometry (triangle mesh) with textures; on-device SfM, MVS, Poisson reconstruction, and UV mapping; output USDZ. Well-suited for editable assets, placement in the room editor, and real-time rendering in RealityKit.

    \item \textbf{3D Gaussian Splatting}: Implicit representation (set of Gaussians); training off-device with differentiable splatting; output \texttt{.ply}. High visual quality and view-dependent effects; no explicit mesh, so geometric editing (e.g.\ boolean operations) is harder. The application uses 3DGS for \emph{viewing} and \emph{comparison} alongside Object Capture, not for replacing furniture in the room editor.
\end{itemize}

By supporting both, the application lets users compare the same object or scene as a mesh (Object Capture) and as a Gaussian splat, and provides a foundation for future extensions (e.g.\ hybrid workflows or on-device 3DGS training as future work).

\chapter{Requirements and System Specification}
\label{ch:requirements}

This chapter specifies the intended behaviour of the application from a \textbf{requirements engineering} perspective: who uses the system, which scenarios must succeed, which quality attributes matter, and on which platform the software runs. Requirements are labelled for traceability to design and tests in Chapter~\ref{ch:implementation}.

\section{Target Users and Interaction Scenarios}

\textbf{Primary user.} The end user is a person redesigning a home or office room who owns a compatible iPhone or iPad and wants to preview furniture changes without permanent purchases.

\textbf{Secondary participants.} A developer extending the codebase benefits from a clean architecture.

\textbf{Core use cases} include: (UC-1) capture a room with RoomPlan and save it; (UC-2) open the saved room in the 3D editor; (UC-3) select a detected furniture instance and assign a replacement USDZ model; (UC-4) adjust position, orientation, and scale with gestures; (UC-5) leave the app and return later with edits intact; (UC-6) capture a standalone object with Object Capture and add it to the furniture library; (UC-7) optionally import a 3DGS \texttt{.ply} and view it alongside or
instead of a mesh for comparison; (UC-8) capture a single photograph with
the in-app camera and receive an interactive 3DGS scene reconstructed by
SHARP on a local server within approximately 2-10 seconds.


\section{Functional Requirements}
\label{sec:func_req}

\begin{description}
    \item[R-1 Room capture.] The system shall run a RoomPlan-based capture session and produce a \texttt{CapturedRoom} containing walls, floors, ceilings where detected, doors, windows, and classified objects with transforms and dimensions.

    \item[R-2 Persistence of scan data.] The system shall serialise each room's \texttt{CapturedRoom} to application-accessible storage (e.g.\ JSON under the Documents directory) keyed by room identifier, and reload it on demand.

    \item[R-3 Editor scene.] The system shall construct a three-dimensional scene in RealityKit representing the captured geometry and furniture placeholders or replacements, without requiring a live AR session for editing.

    \item[R-4 Furniture replacement.] For each detected object selected by the user, the system shall allow binding a USDZ asset (catalog or user-captured) and shall scale and position it to approximate the detected bounding volume.

    \item[R-5 Manipulation.] The user shall translate, rotate, and scale replacement models using on-screen gestures or equivalent controls.

    \item[R-6 Object Capture.] The system shall expose a guided object-scanning flow using Object Capture APIs and shall write resulting USDZ files to a location usable by the furniture library.

    \item[R-7 Gaussian splat viewing.] The system shall parse a standard 3DGS \texttt{.ply} with expected attribute layout and render it with GPU splatting for interactive orbit viewing.

    \item [R-8 Room deletion.] The system shall allow the user to delete any saved room from the room list, with a confirmation step to prevent accidental data loss. Deletion shall remove all associated persistent artefacts (room JSON, replacement records, cached PLY files).

    \item [R-9 SHARP single-image inference.] The system shall provide a capture interface that sends a single JPEG image from the device camera to a user-configured SHARP inference server over the local network, receives a 3DGS \texttt{.ply} file in response, and stores it in the app's Documents directory. The HTTP client shall use a request timeout of at least 360\,s to accommodate first-run Metal shader compilation on the server.

    \item [R-10 Server reachability and configuration.] The system shall perform a health-check ping to the configured SHARP server URL before each inference session and display the server's reachability status (reachable, unreachable, or unknown) with a visual indicator. The user shall be able to configure the server base URL within the app at any time.

\end{description}

\section{Non-Functional Requirements}
\label{sec:nfr_req}

\begin{description}
    \item[N-1 Performance.] On supported hardware, room editing shall remain interactive (target: sustained frame rates comfortable for gesture manipulation; exact FPS depend on scene complexity).

    \item[N-2 Privacy.] Room scanning and object reconstruction shall not require uploading scan data to a third-party server for the Apple API paths described in this thesis.

    \item[N-3 Usability.] The user shall receive guidance during RoomPlan and Object Capture sessions (system-provided or app-provided prompts). Terminology in the UI shall remain consistent with RoomPlan semantics where applicable.

    \item[N-4 Maintainability.] Application logic shall be partitioned into models (e.g.\ \texttt{Room}, replacement records), services (persistence, rendering, catalog), and SwiftUI views, following observable patterns suitable for SwiftUI.

    \item[N-5 Portability within iOS.] The project shall target a documented minimum iOS version matching RoomPlan and Object Capture availability; features degrade gracefully when optional hardware (LiDAR) is absent where Apple allows.

    \item [N-6 Network resilience.] The SHARP inference client shall use a dedicated \texttt{URLSession} instance with a 360\,s request timeout and 600\,s resource timeout, independent of \texttt{URLSession.shared}. If the request fails or times out, the system shall display a descriptive error and return the UI to the ready state without requiring an app restart.

    \item [N-7 Responsiveness.] Operations taking more than 0.5\,s (room processing, SHARP inference, photogrammetric reconstruction) shall display a progress indicator. The SHARP capture overlay shall warn the user that the first inference may take approximately 2\,minutes due to Metal shader compilation on the server, with subsequent runs taking approximately 10\,s.

    \item [N-8 Error handling.] The system shall handle without crashing: (a) corrupt or incompatible \texttt{.ply} import; (b) SHARP server returning a non-200 HTTP response; (c) network timeout on SHARP inference; (d) Object Capture session cancelled with insufficient coverage; (e) attempt to delete a room currently open in the editor.

\end{description}

\section{System Specification}
\label{sec:system_spec}

\textbf{Programming language and UI.} Swift~5, SwiftUI for declarative interfaces, Combine or observation for reactive updates.

\textbf{Graphics and AR.} RealityKit for the room editor; ARKit and RoomPlan for capture; Metal for the optional Gaussian splat renderer (\texttt{MTKView} bridge from SwiftUI).

\textbf{Data formats.} The system utilizes local data formats for the core scan contract, ensuring structural persistence without commercial cloud dependencies. Concretely, \texttt{CapturedRoom} is encoded via JSON through \texttt{Codable}; USDZ serves meshes; PLY serves splats.

\textbf{Storage layout (conceptual).} Room scan files reside under a subdirectory of Documents (e.g.\ \texttt{RoomScans/}), with one file per room identifier; in-memory caches may mirror loaded \texttt{CapturedRoom} instances for performance.


\begin{table}[htbp]
\centering
\caption{Requirement identifiers and short descriptions (traceability).}
\label{tab:req_trace}
\setlength{\tabcolsep}{6pt}
\renewcommand{\arraystretch}{1.25}
\begin{tabular}{lp{0.68\textwidth}}
\toprule
\textbf{ID} & \textbf{Summary} \\
\midrule
\rowcolor{gray!8}
R-1 & RoomPlan capture session $\rightarrow$ \texttt{CapturedRoom} \\
R-2 & Serialise / reload room scan JSON per room identifier \\
\rowcolor{gray!8}
R-3 & RealityKit non-AR editor scene with orbit camera \\
R-4 & USDZ asset binding, scaling, and placement per detected object \\
\rowcolor{gray!8}
R-5 & Gesture-based translate / rotate / scale with floor alignment \\
R-6 & Object Capture guided flow $\rightarrow$ USDZ written to furniture library \\
\rowcolor{gray!8}
R-7 & Binary PLY import and Metal GPU splatting renderer \\
R-8 & Room deletion with confirmation and artefact cleanup \\
\rowcolor{gray!8}
R-9 & SHARP single-image inference via local server; 360\,s timeout \\
R-10 & Server health-check ping; reachability indicator; URL configuration \\
\midrule
\rowcolor{gray!8}
N-1 & Rendering performance: interactive FPS; splat budget $\leq$300\,k Gaussians \\
N-2 & Privacy: no third-party upload; SHARP traffic to LAN server only \\
\rowcolor{gray!8}
N-3 & Usability: guided flows, splash screen, actionable error messages \\
N-4 & Maintainability: MVVM, decoupled services, \texttt{ObservableObject} patterns \\
\rowcolor{gray!8}
N-5 & Portability: iOS~16+ minimum; LiDAR features degrade gracefully \\
N-6 & Network resilience: custom \texttt{URLSession} timeouts; graceful failure recovery \\
\rowcolor{gray!8}
N-7 & Responsiveness: progress indicators for all operations $>$0.5\,s \\
N-8 & Error handling: no crash on PLY corruption, server error, or timeout \\
\bottomrule
\end{tabular}
\end{table}


\begin{figure}[H]
\centering
\thesisfig{0.85\textwidth}{figures/ch4_context_diagram.pdf}
\caption{System context: user, application, and main persistent artefacts (USDZ, PLY, room JSON).}
\label{fig:context_req}
\end{figure}

\section{Constraints and Dependencies}

The application depends on Apple's RoomPlan and Object Capture frameworks and on device capabilities documented by Apple \cite{apple2024roomplan2}. 3DGS viewing depends on compatible \texttt{.ply} exports from external training code or from the SHARP inference path (R-9), which is a \emph{computational} service on the user's network, distinct from third-party content hosting. No requirement in this chapter mandates a specific commercial marketplace; the core thesis pipeline focuses on spatial reconstruction and local editing.

\section{Risk Analysis}

Table~\ref{tab:risks} identifies the principal technical risks and the
mitigation strategies implemented in the application.

\begin{table}[h]
\centering
\caption{Risk analysis: risks, likelihood, impact, and mitigation.}
\label{tab:risks}
\setlength{\tabcolsep}{6pt}
\renewcommand{\arraystretch}{1.3}
\begin{tabular}{p{3.0cm} c c p{5.2cm}}
\toprule
\textbf{Risk} & \textbf{Likelihood} & \textbf{Impact} & \textbf{Mitigation} \\
\midrule
\rowcolor{gray!8}
SHARP server unavailable
  & Medium & Medium
  & Health-check ping before capture; graceful error card in UI; all other
    features remain independent of server. \\
SHARP first-run timeout ($\approx$2\,min)
  & High & Medium
  & Custom \texttt{URLSession} with 360\,s timeout; UI overlay explicitly
    warns user of first-run shader compilation delay. \\
\rowcolor{gray!8}
Rendering lag from large \texttt{.ply}
  & High & High
  & Server-side opacity pruning to ${\leq}$300\,k Gaussians before
    returning response to device. \\
RoomPlan incomplete on non-LiDAR device
  & Medium & High
  & Feature gated behind LiDAR capability check; non-LiDAR devices access
    SHARP and Object Capture only. \\
\rowcolor{gray!8}
Object Capture insufficient coverage
  & Medium & Medium
  & Minimum image count enforced; system guidance prompts user to cover
    all angles. \\
Corrupt \texttt{.ply} import
  & Low & Medium
  & Typed parser error surfaced as alert; recovery suggestion provided. \\
\rowcolor{gray!8}
Accidental room deletion
  & Low & High
  & Confirmation alert with destructive \emph{Delete} label; cancel is
    the default. \\
\bottomrule
\end{tabular}
\end{table}



\chapter{System Architecture, Implementation and Evaluation}
\label{ch:implementation}

This chapter describes how the requirements of Chapter~\ref{ch:requirements} are realised in software. It presents the \textbf{architectural decomposition}, the \textbf{data path} from RoomPlan scan to persisted room and RealityKit entities, the roles of major Swift types and services, the \textbf{user interface} structure, and a \textbf{testing} approach suitable for a bachelor project.

\section{Architectural Overview}

The application entry point is a SwiftUI \texttt{App} type that configures core services and hosts a root view responsible for the local interior-design workflow, including room scanning, editing, and object capture.

Architecturally, the code follows patterns compatible with \textbf{SwiftUI and MVVM}: views observe reference-type \texttt{ObservableObject} managers (e.g.\ \texttt{RoomManager}, \texttt{FurnitureCatalogManager}, \texttt{ScannedFurnitureManager}, \texttt{ProjectManager}) injected through the environment or passed as \texttt{StateObject} / \texttt{ObservedObject}. Long-running capture sessions (RoomPlan, Object Capture) live in dedicated views that bridge UIKit or system controllers where needed. Rendering uses \textbf{RealityKit} (\texttt{ARView} in non-AR camera mode) for the room editor and \textbf{Metal} via \texttt{MTKView} for Gaussian splats.

The architecture is structured into four explicit layers, as shown in Figure~\ref{fig:architecture}:

\begin{enumerate}
    \item \textbf{SwiftUI Views} - declarative UI components such as \texttt{MainTabBar}, \texttt{RoomListView}, \texttt{RoomDesignerView}, \texttt{RoomScanView}, \texttt{ObjectScanView}, \texttt{GaussianSplatView}, \texttt{SHARPCaptureView}, and \texttt{FurnitureLibraryView}.
    \item \textbf{Observable Managers} - \texttt{ObservableObject} classes holding domain state (\texttt{RoomManager}, \texttt{FurnitureCatalogManager}, \texttt{ScannedFurnitureManager}, \texttt{ProjectManager}, \texttt{DecorManager}), injected as environment objects or \texttt{StateObject} references.
    \item \textbf{Services and Renderers} - stateless or singleton-style helpers: \texttt{SHARPService} (HTTP client), \texttt{GaussianSplatRenderer} (Metal GPU pipeline), \texttt{SmartSuggestionsService}, \texttt{RoomColorCaptureService}, and \texttt{HistoryManager}.
    \item \textbf{Apple Frameworks and Storage} - RoomPlan, Object Capture, RealityKit, and Metal at the framework boundary, with persistent artefacts stored in the app's Documents directory.
\end{enumerate}

\begin{figure}[htbp]
\centering
\input{tikz_diagrams/ch5_arch_layers}
% \thesisfig{0.95\textwidth}{figures/ch5_arch_layers.png}
\caption{Layered application architecture: SwiftUI views observe injectable managers, which delegate to services and Apple frameworks and Storage artefacts (room JSON, USDZ, PLY) }
\label{fig:architecture}
\end{figure}

\section{Data Model and Persistence}
\label{sec:ch5_data_model}

\subsection{Room aggregate}

The domain type \texttt{Room} (Swift \texttt{struct}, \texttt{Codable}) aggregates identity (\texttt{id}, \texttt{userId}), timestamps, a list of \texttt{FurnitureReplacement} records, optional \texttt{RoomMaterials} for wall/floor colours captured during scan, optional \texttt{projectId}, and legacy fields kept for backward compatibility. It does \textbf{not} embed the full \texttt{CapturedRoom} inside the JSON that lists all rooms; instead, heavy scan data is stored \textbf{separately}.

\subsection{\texttt{CapturedRoom} storage}

An extension on \texttt{Room} manages \texttt{CapturedRoom} persistence:

\begin{itemize}
    \item \textbf{Directory.} \texttt{Documents/RoomScans/} created on demand.

    \item \textbf{File naming.} One file per room, e.g.\ \texttt{\{roomUUID\}.json}, chosen so that renames or metadata updates do not collide with scan blobs.

    \item \textbf{In-memory cache.} A static dictionary maps room identifiers to loaded \texttt{CapturedRoom} instances to avoid repeated disk I/O while editing.

    \item \textbf{API.} Callers use \texttt{setCapturedRoom(\_:)}, \texttt{getCapturedRoom()}, and helpers that check whether scan data exists (\texttt{hasScanData}).
\end{itemize}

\texttt{RoomManager} owns the array of \texttt{Room} values, serialises it under a user-scoped key (e.g.\ \texttt{saved\_rooms} combined with user identity), and exposes create/update/delete operations. Duplicate detection merges rapid double-submissions of the same room name.

\subsection{Two-level persistence cache}

The persistence layer uses a two-level strategy to balance load latency and memory consumption:

\begin{enumerate}
    \item \textbf{Level 1 - in-memory.}
    \texttt{Room.inMemoryRooms} is a static \texttt{[UUID: CapturedRoom]} dictionary. A room's \texttt{CapturedRoom} loaded from disk during an editor session is stored here and served from memory on subsequent accesses within the same process lifetime, avoiding repeated JSON decoding of potentially large scan files.

    \item \textbf{Level 2 - disk.}
    \texttt{Documents/RoomScans/\{uuid\}.json} provides durable storage. Files are written atomically (\texttt{.atomic} write option) so a crash mid-write does not corrupt the scan.
\end{enumerate}

\texttt{Room.cleanOrphanedScanFiles(existingRoomIds:)} is called during startup to delete JSON files whose UUIDs no longer match any room in the saved-rooms index, preventing unbounded storage growth after rooms are deleted (R-8).

\subsection{Furniture replacement records}

\texttt{FurnitureReplacement} ties a RoomPlan object identifier (or equivalent key) to a USDZ resource name or path, plus transform overrides the editor applies. Persisting these records satisfies R-4 and R-5 together with the scan file.

Beyond direct USDZ replacement of RoomPlan-detected furniture,
\\ 
\texttt{FurnitureReplacement} supports three \emph{virtual decor} categories that are not necessarily tied to a physically detected object:

\begin{description}
    \item[Wall art.] \texttt{WallArtPickerView} lets the user select a frame style (\texttt{FrameStyle}) and size (\texttt{FrameSize}), and optionally provide a photo from the camera roll. The JPEG is saved to a per-app directory and referenced by filename in \\ \texttt{wallArtImageFile}. The RealityKit entity applies the image as a texture on a plane mesh sized to the chosen frame dimensions, anchored to the nearest detected wall surface.

    \item[Carpet.] \texttt{CarpetPickerView} exposes pattern, colour (\texttt{carpetColorRGB}), and size presets plus a custom width/depth in centimetres. The entity is a flat rectangular mesh placed on the floor plane.

    \item[Lighting.] \texttt{LightingPickerView} maps to a RealityKit point or spot light  \\(\texttt{LightFixtureType}), with RGB colour and luminous intensity in lumens. The entity is placed at a user-chosen ceiling or wall position and affects the PBR shading of other entities in the same scene.
\end{description}

All three categories share the same \texttt{FurnitureReplacement} persistence path and re-materialise from stored field values on scene reload, discriminated by the \texttt{originalCategory} string (\texttt{"wallArt"}, \texttt{"carpet"}, \texttt{"lighting"}).

The full domain model, types, fields, and storage relationships, is shown in Figure~\ref{fig:domain_model}.

\begin{figure}[H]
\centering
\input{tikz_diagrams/ch5_domain_model}
% \thesisfig{0.97\textwidth}{figures/ch5_domain_model.png}
\caption{Domain model: \texttt{Room} aggregates \texttt{FurnitureReplacement} records (including decor categories) and optionally references a \texttt{Project}. The heavy \texttt{CapturedRoom} is stored separately on disk, keeping the in-memory room list lightweight.}
\label{fig:domain_model}
\end{figure}

\subsection{On-device storage layout}

Figure~\ref{fig:storage_layout} summarises the full file-system layout. The separation between lightweight metadata in UserDefaults and heavy scan artefacts in the Documents directory is the architectural decision that makes the room list load instantly regardless of the number or size of stored scans.

\begin{figure}[H]
\centering
\input{tikz_diagrams/ch5_storage_layout}
% \thesisfig{0.88\textwidth}{figures/ch5_storage_layout.png}
\caption{On-device storage layout. Binary files (room scan JSON, USDZ assets, wall art images, and SHARP PLY outputs) are segregated into typed subdirectories under the app's Documents folder. Lightweight per-user lists are encoded as JSON and persisted in UserDefaults. Registered user accounts, including password hashes and Google identifiers,  are stored in a SwiftData-managed SQLite database. Authentication tokens are kept separately in the system Keychain.}
\label{fig:storage_layout}
\end{figure}

\section{Data Flow: Scan to Editor}

Here is described the end-to-end data flow from the initial RoomPlan capture through persistence to the interactive editor session.

\begin{enumerate}
    \item \textbf{Capture.} \texttt{RoomCaptureView} / session (RoomPlan) runs until the user finalises; the app receives a \texttt{CapturedRoom}.

    \item \textbf{Commit.} \texttt{RoomManager.createRoom} (or update path) stores the structure: \texttt{Room} list saved, \texttt{setCapturedRoom} writes JSON to \texttt{RoomScans/}.

    \item \textbf{Open editor.} The designer view loads \texttt{getCapturedRoom()}, builds RealityKit entities for walls, floor, openings, and for each object either a placeholder or a resolved USDZ from replacement state.

    \item \textbf{Edit.} Gestures update entity transforms; changes write back to \\  \texttt{FurnitureReplacement} and dirty the \texttt{Room} aggregate; \\ \texttt{RoomManager.updateRoom} persists.

    \item \textbf{Relaunch.} \texttt{RoomManager.loadRooms} restores the catalogue; selecting a room reloads scan JSON from disk into the cache and reconstructs the scene.
\end{enumerate}

% \begin{figure}[H]
% \centering
% \thesisfig{0.97\textwidth}{figures/ch5_dataflow.png}
% \caption{Swimlane data-flow diagram: the scan captured in \texttt{RoomScanView} is handed off to \texttt{RoomManager} for persistence; the designer loads from disk or cache and writes back after each gesture or furniture replacement. The flow is a deliberate hand-off between Apple's capture session and the application's own RealityKit scene graph.}
% \label{fig:dataflow}
% \end{figure}

\section{Key Modules and Views}

\subsection{Navigation and tabs}

\texttt{MainTabBar} provides primary navigation among high-level areas such as \textbf{Furniture}, \textbf{Rooms}, and \textbf{Projects}, using SwiftUI tab buttons styled for the application's theme. This separates catalogue browsing from per-room workflows.

\subsection{Room scanning UI}

Views such as \texttt{RoomScanView} (and related flow coordinators) embed RoomPlan's capture experience, handle session lifecycle, and pass the resulting \texttt{CapturedRoom} to room creation logic. The session is finalised by calling \\ \texttt{RoomCaptureSession.stop(pauseARSession: true)}, which yields a complete \texttt{CapturedRoom} instance containing typed structural geometry, object identifiers, and transforms ready for persistence.

\subsection{Room designer}

\texttt{RoomDesignerView} is the largest single view module: it hosts the RealityKit \texttt{ARView}, configures lighting and non-AR camera behaviour, spawns entities for plan geometry, wires gesture recognisers for selection and manipulation, coordinates furniture picking from \texttt{FurnitureCatalogManager} and scanned assets, and integrates optional features (materials, measurements, suggestions). Replacement logic aligns USDZ model bounds to RoomPlan dimensions using uniform scaling and grounding rules described in Section~\ref{sec:alignment}.

\subsection{Object Capture}

\texttt{ObjectScanView} drives \texttt{ObjectCaptureSession}, displays guidance, collects images, and runs \texttt{PhotogrammetrySession} to produce USDZ. Completed models are registered through \texttt{ScannedFurnitureManager} so the designer can pick them like catalogue items.

\subsection{Gaussian splat pipeline}

\texttt{GaussianSplatData.swift} parses binary or ASCII PLY layouts emitted by common 3DGS trainers, normalising property names (\texttt{f\_dc\_*}, scales, quaternions). \\ \texttt{GaussianSplatRenderer.swift} uploads instances to GPU buffers, sorts by depth each frame, and issues draw calls with Metal shaders that approximate 2D Gaussian falloff and blending. \texttt{GaussianSplatView} wraps an \texttt{MTKView} for SwiftUI presentation.

\subsection{SHARP single-image inference}
\label{sec:sharp_implementation}

SHARP (\emph{Sharp Monocular View Synthesis}~\cite{sharp-monocular-view}) is a monocular novel-view synthesis model that accepts a single JPEG image and returns a 3DGS \texttt{.ply} file in approximately 10\,s on a GPU-equipped Mac, or up to 2\,minutes on the first run while Metal shaders compile on the server.

\texttt{SHARPService} is a pure-static Swift enum that encapsulates the HTTP client contract. Its design addresses three non-trivial constraints:

\begin{enumerate}
    \item \textbf{Timeout.} \texttt{URLSession.shared} caps at 60\,s regardless of \\ \texttt{URLRequest.timeoutInterval}. The service creates a dedicated \texttt{URLSession} with \texttt{timeoutIntervalForRequest} = 360\,s and \texttt{timeoutIntervalForResource} = 600\,s, isolating long inference requests from other app network activity (N-6).

    \item \textbf{Server configuration.} The base URL is stored in \texttt{UserDefaults} under the key \texttt{sharp\_server\_base\_url} and exposed through a text field in the capture UI. This avoids hardcoding a LAN address and lets the user reconfigure without rebuilding (R-10).

    \item \textbf{Health-check.} Before opening the camera, the app issues a \texttt{GET /health} ping on a separate 4\,s-timeout session. A colour-coded indicator informs the user of reachability before they spend time composing the shot (R-10). If the ping fails, the capture button is disabled with an actionable error message.
\end{enumerate}

The inference path (\texttt{POST /sharp/infer}) sends the JPEG body (92\,\% quality) and writes the raw response bytes as \texttt{sharp\_\{timestamp\}.ply} to the Documents directory, immediately openable in \texttt{GaussianSplatView} without any additional processing step. Error variants (\texttt{SHARPError}) cover missing URL, JPEG encoding failure, non-200 HTTP responses, network errors, and disk write failure; all surfaced as localised alert strings without crashing (N-8).

% \begin{figure}[H]
% \centering
% \thesisfig{0.97\textwidth}{figures/ch5_sharp_pipeline.png}
% \caption{SHARP inference pipeline: the device captures a JPEG, posts it to the local SHARP server on the LAN, receives a \texttt{.ply}, and opens it in the Metal splat renderer. A health-check ping with a 4\,s timeout precedes each capture session to provide the reachability indicator.}
% \label{fig:sharp_pipeline}
% \end{figure}

\subsection{Supporting services}

\texttt{DecorManager} owns wall art and carpet asset directories; \texttt{MeasurementService} computes room dimensions from the \texttt{CapturedRoom} geometry for the statistics overlay; \texttt{SmartSuggestionsService} proposes furniture placement hints; \\ \texttt{RoomColorCaptureService} samples ARKit frames during scanning to capture dominant wall and floor colours into \texttt{RoomMaterials}; \texttt{ProjectManager} groups multiple rooms into named projects. These sit alongside the core pipeline and may persist auxiliary JSON or asset references.

\section{USDZ Alignment Heuristics}
\label{sec:alignment}

When a replacement is first bound, the editor must align the USDZ asset to the RoomPlan-detected bounding volume without distorting the model. Two rules are applied:

\begin{enumerate}
    \item \textbf{Uniform scale.} The USDZ model's axis-aligned bounding box is computed in model space. The scale factor $s$ is:
    \[
        s = \min\!\left(\frac{d_x^{\mathrm{room}}}{d_x^{\mathrm{usdz}}},\;
                         \frac{d_y^{\mathrm{room}}}{d_y^{\mathrm{usdz}}},\;
                         \frac{d_z^{\mathrm{room}}}{d_z^{\mathrm{usdz}}}\right)
    \]
    This fits the largest dimension without exceeding the detected bounds on any axis while preserving the model's aspect ratio.

    \item \textbf{Floor grounding.} After scaling, the entity's local $y$-origin is translated so that the bottom face of the scaled model sits flush with the detected object's base plane (derived from the \texttt{transform} field of the \texttt{CapturedRoom.Object}).
\end{enumerate}

The \texttt{customScale} field in \texttt{FurnitureReplacement} stores an additional user-applied multiplier on top of this automatic fit, allowing fine-tuning after initial alignment without discarding the grounding state.

\section{Metal Gaussian Splatting: Rendering Pipeline}
\label{sec:splat_rendering}

The splat renderer follows the tile-based alpha-compositing formulation from Kerbl et al.\ \cite{kerbl20233dgs}, adapted for single-pass instanced rendering on iOS Metal. Figure~\ref{fig:splat_pipeline} shows the per-frame pipeline.

\begin{figure}[H]
\centering
\input{tikz_diagrams/ch5_splat_pipeline}
% \thesisfig{0.55\textwidth}{figures/ch5_splat_pipeline.png}
\caption{Per-frame Metal Gaussian splatting pipeline: orbit camera matrices are computed on CPU, Gaussians are sorted back-to-front by view-space depth, and a single instanced draw call executes the vertex and fragment shaders that project 3D covariance to 2D and evaluate the Gaussian falloff per pixel.}
\label{fig:splat_pipeline}
\end{figure}

\paragraph{3D covariance.}
Each Gaussian stores a unit quaternion $\mathbf{q}$ (converted to rotation matrix $\mathbf{R}$ via the \texttt{quatToRotMat} Metal function) and a scale vector $\mathbf{s}$. The 3D covariance is:
\[
    \boldsymbol{\Sigma} = \mathbf{M}\mathbf{M}^\top,\quad
    \mathbf{M} = \mathbf{R}\,\mathrm{diag}(\mathbf{s})
\]

\paragraph{2D projection.}
Given the Jacobian $\mathbf{J}$ of perspective projection at the Gaussian centre (camera-space position $\mathbf{t}$, focal lengths $f_x, f_y$) and the view-space rotation $\mathbf{W}$, the screen-space covariance is:
\[
    \boldsymbol{\Sigma}' = \mathbf{T}\boldsymbol{\Sigma}\mathbf{T}^\top,\quad
    \mathbf{T} = \mathbf{J}\mathbf{W},\quad
    \mathbf{J} = \begin{pmatrix} f_x/d & 0 & f_x t_x/d^2 \\ 0 & f_y/d & f_y t_y/d^2 \\ 0 & 0 & 0 \end{pmatrix}
\]
where $d = -t_z$ is the positive depth. A low-pass isotropic filter (adding 0.3 to the diagonal of $\boldsymbol{\Sigma}'$) prevents sub-pixel Gaussians from flickering.

\paragraph{Quad sizing.}
The screen-space rendering quad's half-extent is $3\sqrt{\lambda_{\max}}$, where $\lambda_{\max}$ is the dominant eigenvalue of $\boldsymbol{\Sigma}'$, computed analytically from its trace and determinant to avoid a general eigensolver:
\[
    \lambda_{\max} = \frac{\mathrm{tr}}{2} + \sqrt{\left(\frac{\mathrm{tr}}{2}\right)^2 - \det(\boldsymbol{\Sigma}')}
\]

\paragraph{Fragment evaluation.}
Each fragment evaluates the 2D Gaussian exponent:
\[
    \alpha = \sigma \cdot \exp\!\left(-\tfrac{1}{2}\,
        \mathbf{d}^\top {\boldsymbol{\Sigma}'}^{-1} \mathbf{d}\right),\quad
    {\boldsymbol{\Sigma}'}^{-1} = \frac{1}{\det}\begin{pmatrix}c & -b \\ -b & a\end{pmatrix}
\]
where $\mathbf{d}$ is the pixel offset from the projected Gaussian centre, $\sigma$ is the per-Gaussian stored opacity, and $(a,b,c)$ are the elements of the 2D covariance passed from the vertex stage. Fragments with $\alpha < 1/255$ are discarded early; the remaining fragments use premultiplied-alpha blending \\ (\texttt{sourceRGBBlendFactor = .one}, \texttt{destinationRGBBlendFactor = \\ .oneMinusSourceAlpha}).

\paragraph{Depth sorting.}
Per-frame back-to-front ordering is achieved by projecting each Gaussian centre onto the view-space $z$-axis (dot product with the third row of the view matrix) and sorting with Swift's \texttt{Array.sort}, running in $O(n \log n)$ on CPU. For budgets up to 300\,k Gaussians (the server-side pruning limit, N-1), this takes approximately 10–30\,ms on an A-series CPU, remaining within a 60\,Hz frame budget when combined with GPU draw time.

\section{On-Device versus Off-Device Processing}

RoomPlan scan interpretation, Object Capture reconstruction, and USDZ display are \textbf{on-device}. Training 3D Gaussian Splatting from many views is \textbf{off-device} in the shipped workflow: the user (or developer) runs external tooling, then imports \texttt{.ply} for R-7.

The SHARP path (R-9) is a hybrid: the device captures and sends a single image, but inference runs on a user-controlled Mac on the local network, not on a third-party cloud server, preserving the privacy guarantee of N-2.

\section{User Interface Walkthrough}

This section gives an engineering-oriented \textbf{end-to-end user walkthrough} aligned with the implementation: it supports laboratory demonstrations, requirement traceability, and onboarding documentation.

\subsection{End-to-end user walkthrough}

\textbf{Scanning with RoomPlan.}
The user initiates a new room from the rooms list, which presents a RoomPlan-based capture surface (\texttt{RoomCaptureSession}). While the session runs, the framework performs real-time structure-from-motion and semantic analysis: walls, floor, ceiling, openings, and classified furniture instances appear incrementally in the preview, giving immediate feedback that coverage and lighting are sufficient. The user walks the space to reduce occlusions and improve metric consistency. Finalising the scan is an explicit UI action (\emph{Done}); the view layer stops the capture pipeline by invoking \texttt{RoomCaptureSession.stop(pauseARSession: true)}, which yields a complete \texttt{CapturedRoom} instance containing typed geometry, object identifiers, and transforms ready for persistence.

\textbf{Transition from capture to editor.}
After stop, the application does not keep the user inside a live AR viewport for editing. Instead, the \texttt{CapturedRoom} is handed to the persistence layer: \texttt{RoomManager} creates or updates a \texttt{Room} aggregate, serialises lightweight metadata to the saved-rooms index, and writes the heavy scan JSON under \texttt{Documents/RoomScans/} as described in Section~\ref{sec:ch5_data_model}. When the user opens the designer, the view loads the scan from disk (or cache), instantiates a RealityKit \texttt{ARView} configured for \textbf{non-AR} camera mode, and constructs entities for structural meshes and furniture placeholders. The transition is therefore a deliberate hand-off from Apple's capture session to an orbit-style, mesh-centric scene graph under application control.

\textbf{Furniture replacement in the designer.}
Inside \texttt{RoomDesignerView}, the user selects a detected object (tap or equivalent selection affordance). The UI surfaces the furniture catalogue and user-scanned USDZ assets; choosing a model binds it to the RoomPlan object's identifier and applies initial alignment heuristics (uniform scale and grounding to the detected bounds, Section~\ref{sec:alignment}). The on-screen result is a concrete USDZ mesh replacing the procedural placeholder while remaining associated with the same logical furniture instance.

\textbf{Manipulation, gizmos, and \texttt{FurnitureReplacement} updates.}
Fine placement uses two complementary mechanisms. A \textbf{gizmo} (or axis-constrained manipulator) exposes discrete handles for translation along the world $X$, $Y$, and $Z$ axes, enabling millimetre-level adjustments without ambiguity between degrees of freedom. Complementarily, standard multi-touch \textbf{gestures} provide intuitive rotation (two-finger twist or arcball-style interaction) and uniform scaling (pinch). Each change updates the corresponding RealityKit entity transform; the editor mirrors that state into a \texttt{FurnitureReplacement} record (identifier, asset reference, \texttt{positionOffset}, \texttt{yRotation}, \texttt{customScale}) so that the domain model remains the single source of truth for persistence. Saving or navigating back triggers \\ \texttt{RoomManager.updateRoom}, satisfying R-4 and R-5.

\textbf{Decor items: wall art, carpet, and lighting.}
The decor workflow uses the same \texttt{FurnitureReplacement} persistence path but creates virtual entities that are not tied to a physically detected object. Wall art anchors a textured plane to the nearest wall; carpet places a patterned flat mesh on the floor; lighting adds a RealityKit point or spot light with configurable colour and intensity, affecting PBR shading across the scene.

\textbf{Optional Gaussian splat view (computational cloud bridge).}
Independently of the mesh editor, the user may inspect spatial content through \textbf{3D Gaussian Splatting}. When the workflow uses the SHARP inference server (R-9), the device sends a captured image to a user-configured endpoint on the local network (Section~\ref{sec:sharp}); the server performs ML inference and returns a \texttt{.ply} that the application stores and opens in the Metal \texttt{MTKView}-hosted splat renderer. Users may also import offline-trained splats without invoking the server. The splat viewer offers a complementary, view-dependent visualisation that highlights differences between explicit meshes and implicit Gaussian fields, directly fulfilling objective~4 from the introduction.

\subsection{Demonstration checklist}

\textbf{Core flow (required for full traceability):}
\begin{enumerate}
    \item \textbf{Launch and rooms.} Open the app; verify saved rooms and the action to start a new scan.

    \item \textbf{Capture and stop.} Complete a RoomPlan session; confirm \\ \texttt{RoomCaptureSession.stop(pauseARSession: true)} semantics (final \texttt{CapturedRoom}).

    \item \textbf{Editor.} Orbit the scene; replace a detected item with USDZ; adjust with gizmo and gestures; relaunch to confirm R-2 and R-5.

    \item \textbf{Decor.} Add wall art, carpet, and lighting; confirm persistence after relaunch.

    \item \textbf{Room deletion (R-8).} Delete a room; confirm destructive-action alert; verify JSON removed from \texttt{Documents/RoomScans/}.
\end{enumerate}

\textbf{Optional extensions (require additional hardware or server):}
\begin{enumerate}
    \item \textbf{Object Capture.} Complete a small-object scan; confirm USDZ appears in the picker (R-6).

    \item \textbf{SHARP path.} Configure server URL; confirm reachability indicator; capture image; receive and display \texttt{.ply} in splat viewer (R-9, R-10).

    \item \textbf{Splat comparison.} Import a \texttt{.ply} (or use the SHARP output above); orbit and compare visual quality with the mesh editor.
\end{enumerate}

\begin{figure}[H]
\centering
\begin{minipage}[t]{0.47\textwidth}
  \centering
  \thesisfig{\textwidth}{figures/ch5_ui_03_room_scan.png}
  \caption{RoomPlan capture: real-time wall and furniture detection.}
  \label{fig:ui_scan}
\end{minipage}
\hfill
\begin{minipage}[t]{0.47\textwidth}
  \centering
  \thesisfig{\textwidth}{figures/ch5_ui_04_editor.png}
  \caption{Non-AR room editor: orbit camera with USDZ replacements.}
  \label{fig:ui_editor}
\end{minipage}
\caption{App screenshots: RoomPlan capture (left) and the non-AR room editor (right).}
\label{fig:ui_scan_editor}
\end{figure}

\section{Testing}

\subsection{Functional tests}

Manual test cases map requirements to observable outcomes: successful save after scan (R-1, R-2), visible replacement mesh (R-3, R-4), gesture persistence after restart (R-5), completed USDZ from Object Capture (R-6), splat load without crash (R-7), room deletion with artefact cleanup (R-8), SHARP inference and PLY storage (R-9), and server reachability indicator (R-10).

\subsection{Regression and edge cases}

Corrupt or missing JSON should surface a controlled error or empty state; oversized PLY files should be tested for memory and load time; SHARP server timeout (first-run 2-minute delay) should not crash or freeze the UI. Duplicate room creation is mitigated in \texttt{RoomManager} but should be exercised. Attempting to delete a room while it is open in the editor should be handled gracefully (N-8).

\subsection{Performance}

Informal measurement suffices for the thesis: note frame drops when many entities load; optionally record Instruments traces for Metal and RealityKit. The splat renderer should sustain 60\,fps for scenes up to 300\,k Gaussians on supported hardware (N-1).

\begin{table}[htbp]
\centering
\caption{Complete requirement-to-test mapping.}
\label{tab:test_map}
\begin{tabular}{|c|p{0.68\textwidth}|}
\hline
\textbf{Req.} & \textbf{Test procedure} \\
\hline
R-1, R-2 & Scan room; verify \texttt{.json} appears under \texttt{Documents/RoomScans}. \\
\hline
R-3, R-4 & Open editor; confirm walls and at least one object mesh; bind a USDZ. \\
\hline
R-5 & Transform model; kill app; reopen; transforms preserved. \\
\hline
R-6 & Complete object scan; USDZ appears in furniture picker. \\
\hline
R-7 & Import known-good \texttt{.ply}; orbit without artefacts at 60\,fps. \\
\hline
R-8 & Delete a room; confirm alert; verify \texttt{.json} removed from disk; list updated. \\
\hline
R-9 & Configure server URL; capture photo; receive and display \texttt{.ply} in splat viewer. \\
\hline
R-10 & Stop server; reachability indicator turns red; capture button disabled with error. \\
\hline
\end{tabular}
\end{table}

\subsection{Quantitative evaluation: Object Scan}
\label{sec:object_scan_quant_eval}

To complement requirement-oriented functional checks, the Object Scan path (R-6) was evaluated on 72 captures with paired ground-truth dimensions. For each object, the mean absolute dimension error (MAE over width, height, depth), centroid error, and bounding-box IoU were computed from the recorded trials in \texttt{thesis\_data/object\_scan\_evaluation.csv}.

\begin{table}[htbp]
\centering
\caption{Object Scan quantitative summary.}
\label{tab:object_scan_summary}
\small
\begin{tabular}{|l|c|}
\hline
\textbf{Metric} & \textbf{Value} \\
\hline
Number of evaluated captures & 72 \\
\hline
Mean absolute dimension error (mm) & 1.85 \\
\hline
95th percentile dimension error (mm) & 3.72 \\
\hline
Mean centroid error (mm) & 1.57 \\
\hline
Mean bounding-box IoU & 0.814 \\
\hline
\end{tabular}
\end{table}

\begin{table}[htbp]
\centering
\caption{Object Scan results by scene and material conditions.}
\label{tab:object_scan_conditions}
\small
\begin{tabular}{|l|c|c|c|}
\hline
\textbf{Condition} & \textbf{MAE (mm)} & \textbf{Centroid (mm)} & \textbf{IoU} \\
\hline
Plain background & 0.81 & 0.70 & 0.895 \\
\hline
Cluttered background & 2.89 & 2.44 & 0.732 \\
\hline
Opaque objects & 1.37 & 1.19 & 0.840 \\
\hline
Transparent objects & 2.84 & 2.39 & 0.749 \\
\hline
\end{tabular}
\end{table}

Figure~\ref{fig:object_scan_condition_plot} shows the same trend visually: cluttered scenes and transparent materials produce substantially larger geometric errors than plain-background opaque objects. This is consistent with contour ambiguity and unstable silhouettes at object boundaries, especially when transparent surfaces coexist with background texture.

\begin{figure}[htbp]
\centering
\begin{tikzpicture}
\begin{axis}[
    width=0.88\textwidth,
    height=6.0cm,
    ybar,
    ymin=0,
    ymax=3.5,
    bar width=12pt,
    ylabel={Mean dimension MAE (mm)},
    symbolic x coords={bg:plain,bg:cluttered,mat:opaque,mat:transparent},
    xtick=data,
    x tick label style={rotate=15,anchor=east},
    nodes near coords,
    every node near coord/.append style={font=\scriptsize},
    grid=major
]
\addplot[
    fill=blue!35,
    draw=blue!60!black
] table[x=category,y=mean_mae_dim_mm,col sep=comma] {thesis_data/object_scan_summary_groups.csv};
\end{axis}
\end{tikzpicture}
\caption{Object Scan mean absolute dimension error by background and material condition.}
\label{fig:object_scan_condition_plot}
\end{figure}

\subsection{Quantitative evaluation: Room Scan}
\label{sec:room_scan_quant_eval}

Room Scan was evaluated on 36 captures using paired ground-truth room dimensions and window counts from \texttt{thesis\_data/room\_scan\_evaluation.csv}. Across all trials, geometric fidelity remains high: mean absolute error is 1.90\,cm on length, 1.90\,cm on width, and 1.10\,cm on height, with a combined mean of 1.63\,cm.

\begin{table}[htbp]
\centering
\caption{Room Scan dimension fidelity summary.}
\label{tab:room_scan_summary}
\small
\begin{tabular}{|l|c|}
\hline
\textbf{Metric} & \textbf{Value} \\
\hline
Number of evaluated rooms & 36 \\
\hline
Mean absolute length error (cm) & 1.90 \\
\hline
Mean absolute width error (cm) & 1.90 \\
\hline
Mean absolute height error (cm) & 1.10 \\
\hline
Mean absolute dimension error (cm) & 1.63 \\
\hline
95th percentile dimension error (cm) & 2.27 \\
\hline
\end{tabular}
\end{table}

\begin{table}[htbp]
\centering
\caption{Window detection by visibility condition.}
\label{tab:window_visibility}
\small
\begin{tabular}{|l|c|c|c|}
\hline
\textbf{Visibility condition} & \textbf{Trials} & \textbf{Detection success} & \textbf{Missed windows} \\
\hline
Clear view & 22 & 100\% & 0 \\
\hline
Curtains closed & 7 & 0\% & 10 \\
\hline
Partial drape coverage & 7 & 0\% & 11 \\
\hline
\end{tabular}
\end{table}

The results support two key observations for R-1 and R-2. First, metric room reconstruction quality is stable, with small absolute errors across all axes, indicating high dimension fidelity for room envelope estimation. Second, window detection is highly sensitive to visibility: openings are reliably identified when unobstructed, while windows can be missed when covered by curtains or drapes.

\begin{figure}[htbp]
\centering
\begin{tikzpicture}
\begin{axis}[
    width=0.46\textwidth,
    height=5.8cm,
    ybar,
    ymin=0,
    ymax=2.2,
    ylabel={Mean absolute error (cm)},
    symbolic x coords={Length,Width,Height},
    xtick=data,
    nodes near coords,
    every node near coord/.append style={font=\scriptsize},
    bar width=14pt,
    grid=major
]
\addplot[
    fill=teal!40,
    draw=teal!60!black
] table[x=axis,y=mean_error_cm,col sep=comma] {thesis_data/room_dimension_error_summary.csv};
\end{axis}
\end{tikzpicture}
\hfill
\begin{tikzpicture}
\begin{axis}[
    width=0.46\textwidth,
    height=5.8cm,
    ybar,
    ymin=0,
    ymax=110,
    ylabel={Detection success (\%)},
    symbolic x coords={clear,curtains_closed,drapes_partial},
    xticklabels={clear,curtains\_closed,drapes\_partial},
    xtick=data,
    x tick label style={rotate=20,anchor=east},
    nodes near coords,
    every node near coord/.append style={font=\scriptsize},
    bar width=12pt,
    grid=major
]
\addplot[
    fill=orange!55,
    draw=orange!80!black
] table[x=window_visibility,y=success_rate_pct,col sep=comma] {thesis_data/room_window_visibility_summary.csv};
\end{axis}
\end{tikzpicture}
\caption{Room Scan quantitative results: mean per-axis dimension error (left) and window detection success by visibility condition (right).}
\label{fig:room_scan_quant_plots}
\end{figure}

\subsection{Efficient 3DGS experiment design under long render times}
\label{sec:efficient_3dgs_design}

For Asset~B experiments (offline-trained 3DGS), full training and rendering cycles routinely exceed 30\,minutes, so the evaluation protocol follows a minimal defensible design:
\begin{enumerate}
    \item \textbf{Representative subset first.} Candidate scenes are stratified by geometry complexity and texture richness, then a small representative subset is used for early rounds.
    \item \textbf{Staged ablations.} Hyperparameter families (camera count, pruning thresholds, densification interval, and output resolution) are varied in stages instead of full-grid search.
    \item \textbf{Early-stop rule.} Any run whose intermediate validation views stop improving or diverge in artefacts is terminated before full convergence.
    \item \textbf{Confirmatory final pass.} Only the best staged configurations are rerun end-to-end on the full selected set.
\end{enumerate}
This procedure preserves methodological rigor while keeping total experimentation time feasible for a bachelor-scale development cycle.

\section{Summary}

The implementation binds Apple's high-level capture APIs to a persistent, editable RealityKit scene graph through explicit storage of \texttt{CapturedRoom} data and replacement metadata. A two-level persistence cache balances load latency with durable atomic writes. The SHARP integration adds a local-network computational offload path that converts a single photo into a 3DGS scene without third-party cloud dependency. Optional Metal splatting demonstrates a second rendering paradigm — driven by differentiable Gaussian projection rather than mesh rasterisation — without compromising the mesh-based interaction model. Requirements R-1--R-10 and non-functional goals N-1--N-8 are all addressed in the design decisions documented above.


% \chapter{Concluzii}
% %\chapter*{Concluzii}
% \label{conclusions}

% \par Concluzii ...
\chapter{Conclusions}
\label{ch:conclusions}

This thesis investigated how a commodity iOS device can support \textbf{interior design preview} through \textbf{parametric room capture}, \textbf{mesh-based editing}, and optional \textbf{plenoptic visualisation}. The theoretical chapters argued that Apple's Object Capture instantiates a classical \textbf{photogrammetric} chain (features, SfM, MVS-class dense matching, Poisson-class surface reconstruction, texturing), whereas RoomPlan builds on \textbf{depth fusion} and learning to emit a structured \texttt{CapturedRoom}. The second theory chapter positioned \textbf{3D Gaussian Splatting} \cite{kerbl20233dgs} as a real-time differentiable alternative to meshes for \emph{viewing}, not for CAD-style furniture replacement.

\section{Summary of Results}

The software artefact demonstrates a coherent pipeline: \textbf{RoomPlan} $\rightarrow$ serialised scan $\rightarrow$ \textbf{RealityKit} editor $\rightarrow$ USDZ replacements with gestures $\rightarrow$ persistence across sessions. \textbf{Object Capture} extends asset authoring; a \textbf{Metal} splat renderer implements R-7 for imported \texttt{.ply} data. Requirements R-1--R-7 (and associated non-functional goals) are addressed at the design level in Chapter~\ref{ch:requirements} and at the code level in Chapter~\ref{ch:implementation}.

\section{Limitations}

Accuracy is bounded by RoomPlan detection quality, photogrammetry assumptions in Object Capture, and the lack of closed-loop feedback between splats and parametric room semantics. 3DGS training remains external; real-time splat \emph{training} on device is research-grade \cite{Guo26}. Usability has not been quantified with SUS/HARUS in this thesis text; Kandil et al.\ \cite{kandil2021ux} provide a methodology for a future study.

\section{Future Work}

Promising extensions include: formal user evaluation; tighter integration between exported scan images and external COLMAP/3DGS pipelines; level-of-detail for large splats; semantic suggestions informed by work such as K\'an et al.\ \cite{kan2021ar} Mithun et al. \cite{Geo2023} demonstrate that synthetic datasets with rich semantic labels can substantially improve indoor scene understanding for AR applications, a direction complementary to the real-scan approach taken here.

The project shows that \textbf{merging API-level room semantics with an editable mesh workflow} yields a practical middle ground between live AR stickers and professional BIM tools, suitable for iterative interior exploration on mobile hardware.

%\addcontentsline{toc}{chapter}{Concluzii}
%\addcontentsline{toc}{chapter}{Conclusions}

\bibliography{references}

\end{document}
